% Generated draft: requires mathematical and source review.
\documentclass[11pt]{article}
\usepackage[margin=25mm]{geometry}
\usepackage{fontspec}
\setmainfont{Latin Modern Roman}
\usepackage{amsmath,amssymb,mathtools,unicode-math}
\setmathfont{Latin Modern Math}
\usepackage{enumitem,needspace,xurl,hyperref}
\setlength{\parindent}{0pt}
\setlength{\parskip}{5pt}
\setlength{\emergencystretch}{3em}
\newcommand{\N}{\mathbb{N}}
\newcommand{\Z}{\mathbb{Z}}
\newcommand{\Q}{\mathbb{Q}}
\newcommand{\R}{\mathbb{R}}
\newcommand{\C}{\mathbb{C}}
\newcommand{\F}{\mathbb{F}}
\newcommand{\A}{\mathbb{A}}
\newcommand{\PP}{\mathbb{P}}
\newcommand{\E}{\mathbb{E}}
\newcommand{\eps}{\varepsilon}
\newcommand{\dd}{\,\mathrm{d}}
\begin{document}
\section*{UnsolvedMath ID: 701399}
\textbf{Source problem code:} catalogue entry
\subsection{Definitions and mathematical statement}
A group identity is an equation \(w(x_1,\ldots,x_n)=1\), where \(w\) is a group word; a group satisfies this identity if it holds for every substitution of its elements for the variables. A variety of groups is the class of all groups satisfying some set of group identities. It is finitely based if a finite set of identities defines that class, with the group axioms understood.

Let
\[
\mathcal B_4=\{G: g^4=1\text{ for every }g\in G\}.
\]
Here the exponent condition is understood in the usual identity sense: each element has order dividing \(4\), without requiring every group to have exponent exactly \(4\). Groups are not assumed finite.

The question attributed to P. M. Neumann is: for every variety \(\mathcal V\subseteq\mathcal B_4\), does there exist a finite set \(\Sigma\) of group identities such that, for every group \(G\),
\[
G\in\mathcal V
\quad\Longleftrightarrow\quad
G\in\mathcal B_4\ \text{and}\ G\models\Sigma?
\]
Since \(\mathcal B_4\) is defined by the single identity \(x^4=1\), this asks precisely whether every such variety has a finite identity basis among all groups.
\subsection{Short English statement}
Does every variety of groups satisfying \(x^4=1\) admit a finite basis of group identities? The supplied registry labels this question open; that label is not an independently verified assessment.
\subsection{Sources}
\begin{itemize}
\item Local UnsolvedMath registry, permanent ID \(701399\), exact source problem code \texttt{catalogue entry}, title \emph{Kourovka Notebook Problem 17.97}. The supplied record attributes the problem to P. M. Neumann, proposed in 2010, and labels it open.
\item \emph{The Kourovka Notebook (Unsolved Problems in Group Theory)}, 17th issue (2010), Problem 17.97; also identified in the supplied background with the 21st issue edition (2026), \texttt{arXiv:1401.0300v48}. The supplied collection description names E. I. Khukhro and V. D. Mazurov as editors. These bibliographic details are supplied attribution and have not been independently checked.
\end{itemize}
\end{document}
